% GENERATED FILE. Edit canonical lessons, then run npm run generate:books.
% canonical-source-hash: fnv1a64:44c37af9e3e07dc2
% canonical-lessons: RU-C48-house, RU-C48-room, RU-C48-kitchen, RU-C48-window, RU-C48-door

\chapter{House, Room, Kitchen, Window, Door}
\label{ch:ru-house-room-kitchen-window-door}
\hlchaptermodality{48}

\begin{chapteropening}
\textbf{By the end of this chapter:} \emph{I can name a house, a room, a kitchen, a window and a door.}

Complete the last lesson of chapter 48: I can name a house, a room, a kitchen, a window and a door.
\end{chapteropening}

\section[\texorpdfstring{\ru{дом} (dom)}{дом (dom)}]{\ru{дом} (dom) — a house, home}

\label{lesson:RU-C48-house}



\begin{quote}
Before the new one: say the Russian for a mushroom, then the Russian for a glass.
\end{quote}

\subsection*{You'll want to know: \ru{дом}}
\textbf{\ru{дом}} (\emph{dom}) — \textquotedblleft{}a house, home\textquotedblright{}.

\begin{cousinweb}
\textbf{\ru{дом}} $\leftarrow$ Proto-Slavic \emph{*domŭ} — inherited. Vasmer lists Sanskrit \emph{dámas}, Greek \emph{dómos} and Latin \emph{domus}, all \textquotedblleft{}house\textquotedblright{}, from PIE \emph{*dem-}, \textquotedblleft{}to build\textquotedblright{}.
\end{cousinweb}

\begin{grammarlens}[title={What you've built}]
One more word for this chapter.
\end{grammarlens}

\subsection*{Your turn}
\begin{itemize}
\raggedright
  \item \emph{Say it:} \emph{dom}
  \item \emph{Say it:} \emph{dom}, once more
  \item \emph{Say it:} say \emph{dom}
  \item \emph{Recall:} say the Russian for a mushroom, then the Russian for a glass, then say \emph{dom} again
\end{itemize}

\begin{tcolorbox}[breakable,colback=teal!4,colframe=teal!35!black,title={Before you move on}]
What does \ru{дом} mean? (\textquotedblleft{}A house, home\textquotedblright{}.) Say it once more.
\end{tcolorbox}
\section[\texorpdfstring{\ru{комната} (kómnata)}{комната (kómnata)}]{\ru{комната} (kómnata) — a room}

\label{lesson:RU-C48-room}



\begin{quote}
Before the new one: say the Russian for a glass, then the Russian for a house.
\end{quote}

\subsection*{You'll want to know: \ru{комната}}
\textbf{\ru{комната}} (\emph{kómnata}) — \textquotedblleft{}a room\textquotedblright{}.

\begin{cousinweb}
\textbf{\ru{комната}} $\leftarrow$ Latin \emph{(camera) camīnāta}, \textquotedblleft{}a room with a stove\textquotedblright{}, from \emph{camīnus}, \textquotedblleft{}furnace\textquotedblright{}. Vasmer takes the initial stress as a sign it came through Old High German \emph{chemināta}, but the older stress fell on \emph{-\ru{на}-}, which points to Polish, as Kiparsky argued; Jakobson, in Trubachev's note, derived it from Czech.
\end{cousinweb}

\begin{grammarlens}[title={What you've built}]
One more word for this chapter.
\end{grammarlens}

\subsection*{Your turn}
\begin{itemize}
\raggedright
  \item \emph{Say it:} \emph{kómnata}
  \item \emph{Say it:} \emph{kómnata}, once more
  \item \emph{Say it:} say \emph{kómnata}
  \item \emph{Recall:} say the Russian for a glass, then the Russian for a house, then say \emph{kómnata} again
\end{itemize}

\begin{tcolorbox}[breakable,colback=teal!4,colframe=teal!35!black,title={Before you move on}]
What does \ru{комната} mean? (\textquotedblleft{}A room\textquotedblright{}.) Say it once more.
\end{tcolorbox}
\section[\texorpdfstring{\ru{кухня} (kúkhnya)}{кухня (kúkhnya)}]{\ru{кухня} (kúkhnya) — a kitchen}

\label{lesson:RU-C48-kitchen}



\begin{quote}
Before the new one: say the Russian for a house, then the Russian for a room.
\end{quote}

\subsection*{You'll want to know: \ru{кухня}}
\textbf{\ru{кухня}} (\emph{kúkhnya}) — \textquotedblleft{}a kitchen\textquotedblright{}.

\begin{cousinweb}
\textbf{\ru{кухня}} came, from 1717, through Polish \emph{kuchnia} and Czech \emph{kuchyně} from Old High German \emph{kuchina}, from Latin \emph{coquīna}, \textquotedblleft{}kitchen\textquotedblright{}; before that Russians said \emph{\ru{поварня}} (Vasmer).
\end{cousinweb}

\begin{grammarlens}[title={What you've built}]
One more word for this chapter.
\end{grammarlens}

\subsection*{Your turn}
\begin{itemize}
\raggedright
  \item \emph{Say it:} \emph{kúkhnya}
  \item \emph{Say it:} \emph{kúkhnya}, once more
  \item \emph{Say it:} say \emph{kúkhnya}
  \item \emph{Recall:} say the Russian for a house, then the Russian for a room, then say \emph{kúkhnya} again
\end{itemize}

\begin{tcolorbox}[breakable,colback=teal!4,colframe=teal!35!black,title={Before you move on}]
What does \ru{кухня} mean? (\textquotedblleft{}A kitchen\textquotedblright{}.) Say it once more.
\end{tcolorbox}
\section[\texorpdfstring{\ru{окно} (oknó)}{окно (oknó)}]{\ru{окно} (oknó) — a window}

\label{lesson:RU-C48-window}



\begin{quote}
Before the new one: say the Russian for a room, then the Russian for a kitchen.
\end{quote}

\subsection*{You'll want to know: \ru{окно}}
\textbf{\ru{окно}} (\emph{oknó}) — \textquotedblleft{}a window\textquotedblright{}.

\begin{cousinweb}
\textbf{\ru{окно}} $\leftarrow$ Proto-Slavic \emph{*okŭno}, built on \emph{\ru{око}}, \textquotedblleft{}eye\textquotedblright{} (Vasmer): a window is the eye of the house, as English \emph{window} is Old Norse \emph{vindauga}, \textquotedblleft{}wind-eye\textquotedblright{}.
\end{cousinweb}

\begin{grammarlens}[title={What you've built}]
One more word for this chapter.
\end{grammarlens}

\subsection*{Your turn}
\begin{itemize}
\raggedright
  \item \emph{Say it:} \emph{oknó}
  \item \emph{Say it:} \emph{oknó}, once more
  \item \emph{Say it:} say \emph{oknó}
  \item \emph{Recall:} say the Russian for a room, then the Russian for a kitchen, then say \emph{oknó} again
\end{itemize}

\begin{tcolorbox}[breakable,colback=teal!4,colframe=teal!35!black,title={Before you move on}]
What does \ru{окно} mean? (\textquotedblleft{}A window\textquotedblright{}.) Say it once more.
\end{tcolorbox}
\section[\texorpdfstring{\ru{дверь} (dver)}{дверь (dver)}]{\ru{дверь} (dver) — a door}

\label{lesson:RU-C48-door}



\begin{quote}
Before the new one: say the Russian for a kitchen, then the Russian for a window.
\end{quote}

\subsection*{You'll want to know: \ru{дверь}}
\textbf{\ru{дверь}} (\emph{dver}) — \textquotedblleft{}a door\textquotedblright{}.

\begin{cousinweb}
\textbf{\ru{дверь}} $\leftarrow$ Proto-Slavic \emph{*dvĭrĭ} — inherited. Vasmer lists Lithuanian \emph{dùrys}, Gothic \emph{daúr}, Greek \emph{thýra} and Latin \emph{forēs}: English \emph{door} is a cousin, from PIE \emph{*d\textsuperscript{h}wer-}, \textquotedblleft{}door\textquotedblright{}.
\end{cousinweb}

\begin{grammarlens}[title={What you've built}]
One more word for this chapter.
\end{grammarlens}

\subsection*{Your turn}
\begin{itemize}
\raggedright
  \item \emph{Say it:} \emph{dver}
  \item \emph{Say it:} \emph{dver}, once more
  \item \emph{Say it:} say \emph{dver}
  \item \emph{Recall:} say the Russian for a kitchen, then the Russian for a window, then say \emph{dver} again
\end{itemize}

\begin{tcolorbox}[breakable,colback=teal!4,colframe=teal!35!black,title={Before you move on}]
What does \ru{дверь} mean? (\textquotedblleft{}A door\textquotedblright{}.) Say it once more.
\end{tcolorbox}
